%=====================================================================
\chapter{The Machine Learning Workflow and Toolkit}\label{ch:workflow}
%=====================================================================
\locator{1}{Part I --- Foundations}

\begin{outcomes}
By the end of this chapter you will be able to:
\begin{tight}
  \item \textbf{Describe} the stages of a machine learning experiment, from
        framing the problem to reporting the result.
  \item \textbf{Split} data into training, validation and test sets, and
        \textbf{explain} what each one is for.
  \item \textbf{Prepare} features: impute missing values, encode categories and
        scale numbers.
  \item \textbf{Use} the scikit-learn estimator interface and \textbf{combine}
        preparation and model in a pipeline.
  \item \textbf{Recognise} and \textbf{prevent} data leakage, and
        \textbf{establish} a baseline before building any model.
\end{tight}
\end{outcomes}

\begin{prereq}
\chref{introduction} (examples, features, labels, the paradigms) and
\chref{math} (vectors, distance, and why feature scale matters). Basic Python;
Appendix~A is a reference for every library call used in the labs.
\end{prereq}

\begin{keyterms}
workflow $\bullet$ feature types $\bullet$ training set $\bullet$ validation set
$\bullet$ test set $\bullet$ stratified split $\bullet$ imputation $\bullet$
one-hot encoding $\bullet$ standardisation $\bullet$ min--max scaling $\bullet$
estimator $\bullet$ transformer $\bullet$ \texttt{fit} / \texttt{predict} /
\texttt{transform} $\bullet$ pipeline $\bullet$ data leakage $\bullet$ baseline
\end{keyterms}

\section{One Experiment, Start to Finish}

Every chapter from here on teaches a new algorithm, and every lab runs it inside
the same sequence of steps. Learn the sequence once, here, and the algorithms
become interchangeable parts. \dsref{workflow} shows it.

\dsfig{%
\begin{tikzpicture}[font=\scriptsize,
  st/.style={rounded corners=3pt, draw=#1, line width=0.9pt, fill=#1!8,
             text=#1!55!black, font=\scriptsize\bfseries, minimum width=2.75cm,
             minimum height=0.95cm, align=center},
  a/.style={-{Stealth[length=2mm]}, gray!60, line width=0.9pt}]
\node[st=primary]   (s1) at (0,0)    {1. Frame\\{\tiny T, P, E}};
\node[st=primary]   (s2) at (3.1,0)  {2. Collect\\{\tiny and inspect}};
\node[st=accent]    (s3) at (6.2,0)  {3. Split\\{\tiny lock the test set}};
\node[st=secondary] (s4) at (9.3,0)  {4. Prepare\\{\tiny impute, encode, scale}};
\node[st=secondary] (s5) at (12.4,0) {5. Baseline\\{\tiny the number to beat}};
\node[st=greenhl]   (s6) at (12.4,-2.3) {6. Train\\{\tiny on training data}};
\node[st=greenhl]   (s7) at (9.3,-2.3)  {7. Validate\\{\tiny compare, tune}};
\node[st=goldhl]    (s8) at (6.2,-2.3)  {8. Test once\\{\tiny the final number}};
\node[st=goldhl]    (s9) at (3.1,-2.3)  {9. Report\\{\tiny with the baseline}};
\draw[a] (s1) -- (s2); \draw[a] (s2) -- (s3); \draw[a] (s3) -- (s4);
\draw[a] (s4) -- (s5); \draw[a] (s5) -- (s6); \draw[a] (s6) -- (s7);
\draw[a] (s7) -- (s8); \draw[a] (s8) -- (s9);
\draw[a, greenhl!70!black, line width=1pt] (s7.south) .. controls +(0,-0.8) and +(0,-0.8)
     .. node[below, font=\tiny, text=greenhl!45!black] {iterate: new features, new model}
     (s6.south);
\draw[accent, line width=1pt, dashed] (s3.south) -- (s8.north);
\node[font=\tiny, text=accent, anchor=east, align=right] at (6.05,-1.15)
     {test data waits here,\\untouched, until step 8};
\end{tikzpicture}}%
{The workflow every lab follows. The loop between training and validation may run
dozens of times; the test set, split off in step 3, is used exactly once, in step
8.}{workflow}

\begin{conceptbox}[Why the Order Matters]
Steps 3 and 8 are the only two that touch the test set, and they are placed as
far apart as possible on purpose. Everything between them --- every choice of
feature, algorithm and setting --- is made while looking only at training and
validation data. If any of those choices is made while looking at the test set,
the final number measures how well you fitted the test set, not how well the
model will do on the next task, ticket or project.
\end{conceptbox}

\section{Examples, Features and Labels}

A dataset for supervised learning is a table: one row per \term{example}, one
column per \term{feature}, and one column for the \term{label}. Before choosing
any algorithm, decide what type each feature is, because the type decides how
it must be prepared.

\rowcolors{2}{white}{lightbg}
\begin{longtable}{@{}L{3.0cm}L{5.0cm}L{7.0cm}@{}}
\toprule
\rowcolor{primary!15}
\textbf{Feature type} & \textbf{Examples} & \textbf{Preparation} \\
\midrule
\endhead
Numeric, continuous & Budget spent, hours logged, test coverage & Impute gaps;
scale for distance- and gradient-based models \\
Numeric, count & Open bugs, change requests & As above; a log transform often
helps with long tails \\
Categorical, nominal & Team, project type, technology stack & One-hot encode;
never treat the codes as numbers \\
Categorical, ordinal & Priority (low to critical), severity level & Map to ordered
integers, or one-hot if the gaps are not equal \\
Text & Ticket description, commit message & Turn into counts or weights of words
(\chref{bayes}) \\
Identifier & Task ID, ticket number, employee number & Usually drop. An ID that
predicts the label is a leak, not a pattern \\
\bottomrule
\end{longtable}

\section{Splitting the Data}

\begin{definitionbox}[Training, Validation and Test Sets]
The \term{training set} is what the algorithm learns from. The \term{validation
set} is used to compare models and choose settings. The \term{test set} is held
back and used once, at the end, to estimate performance on new data. A typical
split is 60/20/20; with little data, the training and validation roles are
combined by cross-validation (\chref{generalization}).
\end{definitionbox}

\dsfig{%
\begin{tikzpicture}[font=\scriptsize]
\fill[greenhl!18, draw=greenhl] (0,0) rectangle (9.0,0.8);
\fill[secondary!18, draw=secondary] (9.0,0) rectangle (12.0,0.8);
\fill[goldhl!22, draw=goldhl] (12.0,0) rectangle (15.0,0.8);
\node[font=\scriptsize\bfseries, text=greenhl!45!black] at (4.5,0.4) {TRAINING \quad 60\%};
\node[font=\scriptsize\bfseries, text=secondary!55!black] at (10.5,0.4) {VALIDATION 20\%};
\node[font=\scriptsize\bfseries, text=goldhl!45!black] at (13.5,0.4) {TEST 20\%};
\node[font=\scriptsize, text=greenhl!45!black, align=center] at (4.5,-0.45)
     {fit the model's parameters};
\node[font=\scriptsize, text=secondary!55!black, align=center] at (10.5,-0.5)
     {choose between models\\and settings};
\node[font=\scriptsize, text=goldhl!45!black, align=center] at (13.5,-0.5)
     {estimate, once,\\how it will do};
\draw[decorate, decoration={brace, amplitude=4pt, mirror}, gray!60]
     (0,-1.1) -- (11.95,-1.1) node[midway, below=5pt, font=\scriptsize, text=gray!55!black]
     {looked at many times, so its scores are optimistic};
\draw[decorate, decoration={brace, amplitude=4pt, mirror}, gray!60]
     (12.05,-1.1) -- (15.0,-1.1) node[midway, below=5pt, font=\scriptsize, text=gray!55!black]
     {looked at once};
\end{tikzpicture}}%
{The three-way split. Only the test score is an honest estimate of performance on new
data, and it stays honest only if it is looked at once.}{split}

\begin{alertbox}[Split Before You Do Anything Else]
Split first, and stratify: \texttt{train\_test\_split(..., stratify=y)} keeps the
proportion of each class the same in every part, which matters when one class is
rare. Two further rules apply to data with structure. If there are several rows
per project or per developer, keep all of an entity's rows on the same side of
the split. If the data is ordered in time, train on the past and test on the
future --- a random split lets the model peek at tomorrow.
\end{alertbox}

\section{Preparing Features}

\textbf{Missing values.} Most algorithms cannot accept a blank. Either drop the
rows (wasteful, and biased if the gaps are not random) or \term{impute}: fill
each gap with the median of that feature in the training data. When a value is
missing \emph{for a reason} --- a timesheet never filled in during a crunch, a
project that never kept a risk register --- add a yes/no column recording that it
was missing; the absence is itself informative.

\textbf{Categories.} \term{One-hot encoding} replaces a categorical feature with
one 0/1 column per category. \texttt{ptype} $\in \{$Web, Mobile, Data$\}$
becomes three columns, exactly one of which is 1 in each row. Coding the
project types as 1, 2 and 3 instead would tell the model that a Data project is
``three times'' a Web project, which is nonsense that a linear model will happily
use.

\textbf{Scaling.} \chref{math} showed a distance dominated by the feature with the
largest units. \term{Standardisation} rescales each feature to mean 0 and
standard deviation 1:
\[
z = \frac{x - \mu_{\text{train}}}{\sigma_{\text{train}}}
\]
\term{Min--max scaling} maps each feature onto $[0, 1]$ instead:
$x' = (x - \min)/(\max - \min)$. Tree-based models (\chref{trees},
\chref{ensembles}) do not need scaling; nearest neighbours, SVMs, neural networks,
$k$-means, PCA and anything trained by gradient descent do.

\begin{worked}[Standardising with Training Statistics]
Four past releases had test coverage $0.6, 0.8, 0.9, 0.7$ and $2, 4, 6, 8$
open bugs.

\smallskip
\textbf{Coverage:} $\mu = 0.75$;
$\sigma = \sqrt{\tfrac{0.15^2 + 0.05^2 + 0.15^2 + 0.05^2}{4}}
= \sqrt{0.0125} = 0.112$.\\
\textbf{Open bugs:} $\mu = 5$;
$\sigma = \sqrt{\tfrac{9 + 1 + 1 + 9}{4}} = \sqrt{5} = 2.236$.

\smallskip
A new release arrives with coverage $0.5$ and 10 open bugs:
\[
z_{\text{cov}} = \frac{0.5 - 0.75}{0.112} = -2.24, \qquad
z_{\text{bugs}} = \frac{10 - 5}{2.236} = +2.24
\]
\textbf{Both are now equally unusual} --- 2.24 standard deviations from the
training mean, in opposite directions --- where before the bug difference (5)
was twenty times the coverage difference (0.25). Note that $\mu$ and $\sigma$
came from the \emph{training} releases only. The new release is transformed with
the training statistics, never with statistics that include the new release.
\end{worked}

\section{The scikit-learn Interface}

Every algorithm in this course is available in \texttt{scikit-learn}, and every
one of them is used in the same way. That uniformity is the reason the labs can
swap a decision tree for an SVM by changing one line.

\begin{definitionbox}[Estimator and Transformer]
An \term{estimator} learns from data with \texttt{fit(X, y)} and then makes
predictions with \texttt{predict(X)} (and, for classifiers,
\texttt{predict\_proba(X)} for probabilities). A \term{transformer} learns with
\texttt{fit(X)} and changes data with \texttt{transform(X)}; a scaler learns
$\mu$ and $\sigma$ in \texttt{fit} and applies them in \texttt{transform}. Anything
learned during \texttt{fit} is stored in attributes ending in an underscore, such
as \texttt{coef\_} or \texttt{mean\_}.
\end{definitionbox}

\rowcolors{2}{white}{lightbg}
\begin{longtable}{@{}L{5.9cm}L{1.9cm}L{7.2cm}@{}}
\toprule
\rowcolor{primary!15}
\textbf{Class} & \textbf{Chapter} & \textbf{What it is} \\
\midrule
\endhead
\texttt{LinearRegression}, \texttt{Ridge}, \texttt{Lasso} & 5, 7 & Least-squares
regression, with and without a penalty \\
\texttt{LogisticRegression} & 6 & Linear classifier with probabilities \\
\texttt{DecisionTreeClassifier} & 9 & ID3/CART-style tree \\
\texttt{GaussianNB}, \texttt{MultinomialNB} & 10 & Naive Bayes \\
\texttt{KNeighborsClassifier} & 11 & $k$-nearest neighbours \\
\texttt{SVC} & 12 & Support vector machine \\
\texttt{MLPClassifier} & 13 & Multilayer neural network \\
\texttt{RandomForestClassifier},\newline \texttt{AdaBoostClassifier} & 14 & Ensembles \\
\texttt{KMeans},\newline \texttt{AgglomerativeClustering} & 15 & Clustering \\
\texttt{PCA} & 16 & Principal component analysis \\
\texttt{GaussianMixture},\newline \texttt{SelfTrainingClassifier} & 17 & EM mixtures;
semi-supervised wrapper \\
\bottomrule
\end{longtable}

\section{Pipelines and Data Leakage}

\begin{definitionbox}[Data Leakage]
\term{Data leakage} is any route by which information that will not be available
at prediction time reaches the model during training or evaluation. Leakage
makes evaluation scores optimistic, sometimes spectacularly so, and the model
then fails in use.
\end{definitionbox}

Leakage comes in two forms. \term{Target leakage} is a feature that is recorded
after, or because of, the outcome: \texttt{days\_overdue} in a model of which
tasks will finish late. \term{Train--test contamination} is subtler:
preparing the data using statistics that include the test rows. Fit a scaler on
the full dataset and then split, and the training data has been shifted by the
test data's mean --- the model has seen a summary of the exam.

\dsfig{%
\begin{tikzpicture}[font=\scriptsize,
  bx/.style={rounded corners=2pt, draw=#1, fill=#1!10, text=#1!50!black,
             font=\scriptsize\bfseries, minimum width=2.3cm, minimum height=0.7cm,
             align=center},
  a/.style={-{Stealth[length=2mm]}, gray!60, line width=0.9pt}]
\node[font=\small\bfseries, text=accent, anchor=west] at (-1.2,1.0) {Leaky};
\node[bx=neutral] (a1) at (0,0) {All data};
\node[bx=accent] (a2) at (3.0,0) {Fit scaler\\on all rows};
\node[bx=neutral] (a3) at (6.0,0) {Split};
\node[bx=greenhl] (a4) at (9.0,0.45) {Train};
\node[bx=goldhl] (a5) at (9.0,-0.45) {Test};
\draw[a] (a1) -- (a2); \draw[a] (a2) -- (a3); \draw[a] (a3) -- (a4.west);
\draw[a] (a3) -- (a5.west);
\draw[accent, -{Stealth[length=2mm]}, dashed, line width=0.9pt]
     (a5.south) .. controls +(0,-0.9) and +(0,-0.9) .. node[below, font=\tiny]
     {test rows shaped $\mu$ and $\sigma$} (a2.south);
\node[font=\small\bfseries, text=greenhl!55!black, anchor=west] at (-1.2,-2.3) {Correct};
\node[bx=neutral] (b1) at (0,-3.3) {All data};
\node[bx=neutral] (b2) at (3.0,-3.3) {Split};
\node[bx=greenhl] (b3) at (6.0,-2.85) {Train};
\node[bx=goldhl] (b4) at (6.0,-3.75) {Test};
\node[bx=secondary] (b5) at (9.0,-2.85) {Fit scaler\\on train only};
\node[bx=secondary] (b6) at (9.0,-3.75) {Apply same\\$\mu$, $\sigma$};
\draw[a] (b1) -- (b2); \draw[a] (b2) -- (b3.west); \draw[a] (b2) -- (b4.west);
\draw[a] (b3) -- (b5); \draw[a] (b4) -- (b6); \draw[a] (b5) -- (b6);
\node[anchor=north west, align=left, text width=4.2cm, font=\tiny, text=gray!55!black]
     at (10.6,0.9)
     {A \texttt{Pipeline} makes the correct order automatic: calling
      \texttt{fit} on the pipeline fits every step on the training data only, and
      \texttt{predict} applies the stored statistics to new data.};
\end{tikzpicture}}%
{Train--test contamination. Fitting any preparation step before the split lets the test
rows influence training; a pipeline fitted on the training data alone cannot make this
mistake.}{leakage}

\begin{definitionbox}[Pipeline]
A scikit-learn \term{pipeline} chains transformers and a final estimator into one
object with a single \texttt{fit} and \texttt{predict}. Because every step is fitted
only when the pipeline is fitted, and only on the data passed to it, a pipeline
used inside a split or a cross-validation loop cannot leak. A
\texttt{ColumnTransformer} applies different preparation to different columns
--- scaling to the numbers, one-hot encoding to the categories.
\end{definitionbox}

\section{Baselines}

\begin{definitionbox}[Baseline]
A \term{baseline} is the score of the simplest possible predictor: always predict
the most common class (classification) or always predict the training mean
(regression). Every model must be compared with it, and a model that does not
beat it clearly has learned nothing useful.
\end{definitionbox}

Baselines are humbling. On the late-task data of \chref{introduction}'s lab,
``always on time'' already scores 66\% --- more than the hand-written rule. On
commit data where 99\% of commits are clean, ``always clean'' scores 99\% ---
which is why a defect predictor's accuracy means nothing until it is set beside
the baseline (\chref{evaluation}).

\begin{pitfall}
\begin{tight}
  \item \textbf{Preparing before splitting.} Imputing, scaling or selecting
        features on the full dataset leaks the test set into training. Put every
        step in a pipeline.
  \item \textbf{Using the test set to choose.} Every look at the test set that
        changes a decision turns it into a validation set. Use it once.
  \item \textbf{Label-encoding a nominal feature.} Project-type codes 1, 2, 3 impose
        an order that does not exist. One-hot encode.
  \item \textbf{Keeping an identifier.} If \texttt{task\_id} improves accuracy,
        IDs were assigned in a way that encodes the outcome --- a leak.
  \item \textbf{Reporting a score without the baseline.} ``92\% accurate'' is
        meaningless until you know that ``always no'' scores 90\%.
\end{tight}
\end{pitfall}

%---------------------------------------------------------------------
\begin{lab}[A Leak-Proof Pipeline]
Six hundred synthetic past projects with a missing-value problem and one
categorical feature; the label is whether the project finished late. The whole
workflow, in the order of \dsref{workflow}.
\begin{lstlisting}[language=Python]
import numpy as np
import pandas as pd
from sklearn.model_selection import train_test_split
from sklearn.compose import ColumnTransformer
from sklearn.pipeline import Pipeline, make_pipeline
from sklearn.impute import SimpleImputer
from sklearn.preprocessing import StandardScaler, OneHotEncoder
from sklearn.linear_model import LogisticRegression
from sklearn.dummy import DummyClassifier

# --- 600 past projects, with gaps ---------------------------------------
rng = np.random.default_rng(0)
n = 600
df = pd.DataFrame({
    "clarity": rng.uniform(0.3, 1.0, n),         # requirement clarity
    "coverage": rng.uniform(0.0, 1.0, n),        # test coverage
    "changes": rng.poisson(5, n).astype(float),  # change requests
    "ptype": rng.choice(["Web", "Mobile", "Data"], n),
})
risk = 2.6 - 4 * df.clarity - 2 * df.coverage + 0.1 * df.changes
df["late"] = (risk + rng.normal(0, 0.5, n) > 0).astype(int)
df.loc[rng.random(n) < 0.1, "coverage"] = np.nan        # 10% missing
print("late rate:", df.late.mean().round(3))

X, y = df.drop(columns="late"), df["late"]

# --- 1. split FIRST; the test set is locked away ------------------------
X_tr, X_te, y_tr, y_te = train_test_split(
    X, y, test_size=0.2, stratify=y, random_state=0)

# --- 2. every preparation step lives inside the pipeline ----------------
num = ["clarity", "coverage", "changes"]
prep = ColumnTransformer([
    ("num", make_pipeline(SimpleImputer(strategy="median"),
                          StandardScaler()), num),
    ("cat", OneHotEncoder(handle_unknown="ignore"), ["ptype"]),
])
model = Pipeline([("prep", prep), ("clf", LogisticRegression())])

# --- 3. a baseline first, then the model --------------------------------
base = DummyClassifier(strategy="most_frequent").fit(X_tr, y_tr)
print("baseline accuracy:", round(base.score(X_te, y_te), 3))
model.fit(X_tr, y_tr)
print("model accuracy   :", round(model.score(X_te, y_te), 3))

# --- 4. look inside: the scaler learned from training rows only ---------
scaler = model.named_steps["prep"].named_transformers_["num"][-1]
print("scaler means:", scaler.mean_.round(2))
\end{lstlisting}
The baseline scores about 0.68 and the model about 0.84. Now swap the model:
import \texttt{DecisionTreeClassifier} from \texttt{sklearn.tree} and put it where
\texttt{LogisticRegression()} was. Nothing else changes --- that is the point of
the interface.
\end{lab}

%---------------------------------------------------------------------
\begin{checkpoint}
\begin{tightnum}
  \item Why must the scaler's mean and standard deviation be computed from the
        training set only?
  \item A dataset has a column \texttt{method} with values Scrum, Kanban and
        Waterfall. How should it be encoded, and why not as 1, 2, 3?
  \item A model predicting which projects will be cancelled scores 97\% accuracy.
        What single number do you ask for before you are impressed?
\end{tightnum}
\end{checkpoint}

%---------------------------------------------------------------------
\begin{chaptersummary}
\begin{tight}
  \item Every experiment follows the same workflow: frame, collect, split,
        prepare, baseline, train, validate, test once, report.
  \item The training set fits the model, the validation set chooses between
        models, and the test set --- used once --- estimates performance on new
        data.
  \item Impute gaps (and record that they were gaps), one-hot encode nominal
        categories, and scale numbers for any method based on distance or
        gradients.
  \item Every scikit-learn estimator has \texttt{fit} and \texttt{predict}; every
        transformer has \texttt{fit} and \texttt{transform}. Pipelines chain them.
  \item Leakage --- through a feature recorded after the outcome or through
        preparation fitted on test rows --- produces scores the model cannot
        reproduce in use. Pipelines prevent the second kind.
  \item Always report the baseline beside the model.
\end{tight}
\end{chaptersummary}

\begin{reviewq}
  \item \textbf{(MCQ)} Which set is used to choose between a decision tree and an
        SVM? (a)~training (b)~validation (c)~test (d)~all three.
  \item \textbf{(MCQ)} Which model does \emph{not} need its features scaled?
        (a)~$k$-nearest neighbours (b)~SVM (c)~decision tree (d)~$k$-means.
  \item \textbf{(MCQ)} A column recording the date a project was formally closed
        as failed, in a model predicting project failure, is an example of:
        (a)~target leakage (b)~stratification (c)~imputation (d)~a baseline.
  \item \textbf{(Short)} Explain the difference between \texttt{fit},
        \texttt{transform} and \texttt{predict}.
  \item \textbf{(Short)} Why should a time-ordered dataset not be split at random?
  \item \textbf{(Short)} When is it worth adding a ``value was missing'' column
        alongside an imputed feature?
  \item \textbf{(Applied)} Training values of a feature are $10, 20, 30, 40$. Compute
        the standardised value of a new observation of $45$, showing $\mu$ and
        $\sigma$. Then compute its min--max scaled value, and say what is unusual
        about the result.
  \item \textbf{(Applied)} A colleague scales the whole dataset, selects the ten
        features most correlated with the label, then runs 5-fold
        cross-validation and reports 94\%. Identify every leak in this procedure
        and rewrite it so that it cannot leak.
\end{reviewq}
